\documentclass[aspectratio=169]{beamer}
\usetheme{Madrid}
\usecolortheme{seahorse}

\title{Camera Tracking and Virtual Avatars}
\author{Project Overview}
\date{\today}

\begin{document}

\begin{frame}
  \titlepage
\end{frame}

\begin{frame}{Project Summary}
  \begin{itemize}
    \item Browser-based, real-time camera tracking demos for interactive avatars and gameplay.
    \item Five self-contained HTML demos: \textbf{Bear}, \textbf{Frog}, \textbf{Tiger}, \textbf{Hand}, and \textbf{Ping}.
    \item Focus on accessibility: no installation required for online usage.
  \end{itemize}
\end{frame}

\begin{frame}{Core Technology Stack}
  \begin{itemize}
    \item \textbf{MediaPipe} for hand and face landmark detection.
    \item \textbf{HTML5 getUserMedia} for webcam stream acquisition.
    \item \textbf{Canvas \& SVG} for 2D overlays and avatar animation.
    \item \textbf{Three.js} for 3D rendering in the Hand demo.
  \end{itemize}
\end{frame}

\begin{frame}{Shared Processing Pipeline}
  \begin{enumerate}
    \item Capture camera frames from browser webcam API.
    \item Detect landmarks using MediaPipe models.
    \item Derive high-level features (finger fold, eye/mouth openness, hand openness).
    \item Map features to avatar/game controls with smoothing (LERP).
    \item Render visual feedback in SVG, Canvas, or WebGL.
  \end{enumerate}
\end{frame}

\begin{frame}{Demo 1: Bear}
  \begin{itemize}
    \item Tracks one hand to animate five bear body parts.
    \item Finger fold detection compares wrist-to-tip distance against wrist-to-joint distance.
    \item Thumb fold uses a dedicated opposable-thumb heuristic.
    \item Animated via per-part SVG group transformations in an animation loop.
  \end{itemize}
\end{frame}

\begin{frame}{Demo 2: Frog}
  \begin{itemize}
    \item Face landmarks drive frog eyelids and mouth.
    \item Eye and mouth openness are normalized against face scale for robustness.
    \item Temporal smoothing with LERP reduces jitter and improves visual stability.
  \end{itemize}
\end{frame}

\begin{frame}{Demo 3: Tiger}
  \begin{itemize}
    \item Combines face and two-hand tracking in one scene.
    \item Hand positions and angles drive paw transforms.
    \item Face blendshapes and landmarks drive mouth/head expressions.
    \item Uses a draggable picture-in-picture tracking viewport.
  \end{itemize}
\end{frame}

\begin{frame}{Demo 4: Hand}
  \begin{itemize}
    \item Renders a stylized 3D hand skeleton and mesh in Three.js.
    \item MediaPipe 3D hand landmarks are mapped into scene coordinates.
    \item Animal-emoji stickers are attached to finger tips with orientation updates.
  \end{itemize}
\end{frame}

\begin{frame}{Demo 5: Ping}
  \begin{itemize}
    \item Two-player ping pong controlled by left and right hands.
    \item Hand openness score is computed from fingertip-to-wrist aggregate distance.
    \item Score is mapped to paddle Y position with smoothing for stable gameplay.
  \end{itemize}
\end{frame}

\begin{frame}{Outcome and Future Work}
  \begin{itemize}
    \item Demonstrates practical real-time body tracking on common devices.
    \item Delivers engaging, low-friction avatar and game interactions in-browser.
    \item Future directions:
    \begin{itemize}
      \item More robust multi-user tracking.
      \item Improved temporal filtering and latency optimization.
      \item Expanded avatar rigging and richer interaction design.
    \end{itemize}
  \end{itemize}
\end{frame}

\begin{frame}{Repository Demo Links}
  \small
  \begin{itemize}
    \item Bear: \texttt{bear.html}
    \item Frog: \texttt{frog.html}
    \item Hand: \texttt{hand.html}
    \item Tiger: \texttt{tiger.html}
    \item Ping: \texttt{ping.html}
  \end{itemize}
\end{frame}

\end{document}
